\documentclass[12pt, a4paper]{article}
\usepackage[margin=2.5cm]{geometry}
\usepackage{booktabs}
\usepackage{enumitem}
\usepackage{parskip}

\title{\textbf{R25 Sprint --- Business Overview}}
\author{DXC Service Desk Analytics}
\date{Sprint R25 \quad|\quad 08 March -- 04 April 2026}

\begin{document}
\maketitle
\tableofcontents
\newpage

% ─────────────────────────────────────────────────────────────────────────────
\section{What is a Sprint?}

A sprint is a fixed time period during which the team commits to delivering
a defined set of work. For R25, the sprint runs from \textbf{08 March to
04 April 2026} (28 working and calendar days).

At the start of the sprint, the team agrees on a list of tickets to complete.
This list is called the \textbf{committed scope}. The burndown chart then
tracks, day by day, how much of that scope has been delivered and how much
remains.

% ─────────────────────────────────────────────────────────────────────────────
\section{The Two Input Files}

\subsection{Extract.xlsx --- The Live Jira Export}

This file is exported directly from Jira and contains every ticket across all
projects. It is the source of truth for the current state of each issue:
its status, who it is assigned to, when it was resolved, and so on.

It is uploaded to the database regularly so that the dashboards always
reflect the latest situation.

Key information found in this file:

\medskip
\begin{center}
\begin{tabular}{p{4.5cm} p{9cm}}
\toprule
\textbf{Column} & \textbf{What it means} \\
\midrule
\textbf{Key} &
  Unique identifier of the ticket (e.g.\ \texttt{IRPAUTO-18734}).
  Used to link tickets across all files. \\[6pt]

\textbf{Summary} &
  Short description of the work to be done. \\[6pt]

\textbf{Issue Type} &
  Category of work: \emph{Incident}, \emph{Enhancement},
  \emph{Service Request}, \emph{Change}, etc. \\[6pt]

\textbf{Status} &
  Where the ticket currently sits in the workflow
  (see Section~\ref{sec:workflow}). \\[6pt]

\textbf{Priority} &
  Urgency level: Critical, High, Medium, or Low. \\[6pt]

\textbf{Assignee} &
  The team member responsible for resolving the ticket. \\[6pt]

\textbf{Story Points} &
  An estimate of the effort required, agreed during sprint planning.
  Higher = more complex work. \\[6pt]

\textbf{Created} &
  Date and time the ticket was opened in Jira. \\[6pt]

\textbf{Qualification Date} &
  Date the ticket passed the QA gate --- it has been tested and
  validated before being sent to the client. \\[6pt]

\textbf{Resolved} &
  Date Jira recorded a resolution for the ticket. \\[6pt]

\textbf{Closure Date} &
  Date the ticket was formally closed after all steps were completed. \\[6pt]

\textbf{Environment Type} &
  Whether the issue occurred in \emph{Production} or a
  non-production environment (Recette, UAT, etc.). \\[6pt]

\textbf{Project} &
  The client or product the ticket belongs to
  (e.g.\ Cardif, Federale, Argenta). \\

\bottomrule
\end{tabular}
\end{center}

\subsection{Release\_lifecycle\_R25.xlsx --- The Sprint Planning File}

This file is maintained by the project management team and contains two
important sheets.

\subsubsection*{Sheet 1 --- Scope UPDATE}

Lists the tickets that have been \textbf{officially committed} for sprint R25.
Not every ticket in Jira is included --- only those the team has agreed to
deliver during this sprint.

For each committed ticket it records:
\begin{itemize}
  \item The Jira key.
  \item The story point estimate from planning (may differ from Jira if updated later).
  \item The expected sprint and delivery status from the planning perspective.
\end{itemize}

\subsubsection*{Sheet 2 --- Team Availability}

Lists every team member involved in R25 and their \textbf{productive capacity}
for the sprint, expressed in person-days (\emph{Jours-Homme}, JH).

The capacity already accounts for deductions such as vacations, training
days, and transverse activities. The resulting figure (\texttt{cap\_prod})
represents the net time available for delivery work.

% ─────────────────────────────────────────────────────────────────────────────
\section{The Ticket Workflow}
\label{sec:workflow}

A ticket moves through a series of statuses from the moment it is created
until it is closed. The typical journey looks like this:

\medskip
\begin{center}
\begin{tabular}{p{4cm} p{9cm}}
\toprule
\textbf{Status} & \textbf{What it means} \\
\midrule
\textbf{Open / New} &
  Ticket has just been created, no one has started working on it yet. \\[6pt]

\textbf{Analysis / In Progress} &
  A team member is actively investigating or implementing the fix. \\[6pt]

\textbf{Ready for QA} &
  Development or fix is done. The ticket is waiting for quality
  assurance testing. \\[6pt]

\textbf{Ready for Acceptance} &
  QA passed. The ticket is waiting for client validation. \\[6pt]

\textbf{Resolved} &
  The work has been completed and validated. Jira has recorded a
  resolution date. \\[6pt]

\textbf{Closed} &
  All steps are done and the ticket is formally archived.
  This is the final state. \\[6pt]

\textbf{Cancelled} &
  The ticket was abandoned --- the work is no longer needed. \\

\bottomrule
\end{tabular}
\end{center}

\medskip
For the burndown, a ticket counts as \textbf{delivered} once it reaches
\emph{Resolved}, \emph{Closed}, or \emph{Cancelled} --- or once it has a
valid completion date recorded (qualification, resolution, or closure)
during the sprint window.

% ─────────────────────────────────────────────────────────────────────────────
\section{What the Burndown Chart Shows}

The burndown chart answers one question: \textbf{are we going to finish
everything we committed to by the end of the sprint?}

\begin{itemize}
  \item The \textbf{grey dashed line} (ideal) shows what perfect, steady
        progress would look like --- delivering an equal share of the scope
        every single day.

  \item The \textbf{purple line} (actual) shows what was really delivered
        each day, based on the live Jira extract.

  \item When the purple line is \textbf{above} the grey line, the team is
        behind schedule.

  \item When the purple line is \textbf{below} the grey line, the team is
        ahead of schedule.

  \item \textbf{Flat segments} mean no tickets were delivered that day
        (weekend, public holiday, or blocked work).

  \item \textbf{Steep drops} mean several tickets were delivered on the
        same day.
\end{itemize}

The chart can be read in two ways:
\begin{itemize}
  \item \textbf{Story Points} --- measures effort. Closing a large,
        complex ticket counts more than closing a small one.
  \item \textbf{Ticket Count} --- measures volume. Every closed ticket
        counts equally, regardless of size.
\end{itemize}

Comparing both gives a clearer picture: if ticket count looks good but
story points are lagging, the team may be closing easy tickets while
avoiding the heavy ones.

% ─────────────────────────────────────────────────────────────────────────────
\section{Team Capacity and Productivity}

Each team member has a \textbf{productive capacity} (\texttt{cap\_prod})
expressed in person-days for the sprint. This is calculated from the
planning file after deducting:

\begin{itemize}
  \item Vacation and public holidays.
  \item Training and formation days.
  \item Transverse activities (meetings, support, etc.).
  \item A contingency factor applied by management.
\end{itemize}

Team members with a capacity of zero (\texttt{cap\_prod = 0}) are still
tracked for story points delivered --- they are listed separately as
\textbf{Unplanned Delivery}, meaning they contributed work that was not
accounted for in the sprint plan.

% ─────────────────────────────────────────────────────────────────────────────
\section{How the Data Flows End to End}

\begin{enumerate}[leftmargin=*, label=\textbf{Step \arabic{enumi}.}]

  \item \textbf{Planning.}
        The project team fills in the Release Lifecycle file, listing the
        committed tickets and team availability for the sprint.

  \item \textbf{Upload to database.}
        Both files are uploaded via the dashboard. The planning data goes
        into \texttt{r25\_scope} and \texttt{r25\_team\_avail}. The Jira
        extract goes into the main \texttt{issues} table.

  \item \textbf{Daily work.}
        The team resolves tickets in Jira throughout the sprint.

  \item \textbf{Refresh.}
        A new Jira extract is uploaded regularly. The burndown table
        (\texttt{r25\_burndown}) is refreshed with one click, pulling the
        latest resolution dates and recomputing all daily metrics.

  \item \textbf{Reporting.}
        The Streamlit dashboard and Power BI both read from the same
        database tables, ensuring everyone sees the same numbers.

\end{enumerate}

\end{document}
